% =====================================================================
% 03-taxonomy.tex  (core contribution)
% =====================================================================
\section{A Three-Axis Taxonomy}\label{sec:taxonomy}

The diversity of RL methods for LLMs can be overwhelming, but almost
every method can be characterized by answering three independent
questions: \emph{What signal does it learn from? How is the policy
updated? And over what temporal scope?} We therefore organize the field
along three nearly orthogonal axes, illustrated in
Fig.~\ref{fig:taxonomy}.

\begin{figure*}[t]
\centering
\begin{forest}
forked edges,
for tree={grow'=0,draw,rounded corners,font=\footnotesize,align=left,
          edge+={semithick},parent anchor=east,child anchor=west,
          l sep=7mm,s sep=1.2mm,inner sep=2.5pt,}
[{\textbf{RL for}\\\textbf{LLM Training}}, fill=rootgray!18, align=center
  [{\textbf{Axis A}\\Reward Signal}, fill=axisA!18, align=center
    [{Human preference (RLHF)~\cite{ouyang2022instructgpt}}]
    [{AI preference (RLAIF, Constitutional AI)~\cite{bai2022constitutional}}]
    [{Implicit / data preference (DPO)~\cite{rafailov2023dpo}}]
    [{Verifiable / rule-based (RLVR)~\cite{shao2024deepseekmath}}]
    [{Granularity: outcome (ORM) vs.\ process (PRM)~\cite{lightman2023verify}}]
  ]
  [{\textbf{Axis B}\\Optimization}, fill=axisB!18, align=center
    [{Policy gradient + critic (PPO)~\cite{schulman2017ppo}}]
    [{Critic-free / group (REINFORCE, RLOO, GRPO)~\cite{shao2024deepseekmath}}]
    [{Direct / offline (DPO, IPO, KTO)~\cite{rafailov2023dpo}}]
  ]
  [{\textbf{Axis C}\\Training Scope}, fill=axisC!18, align=center
    [{Single-step / single-turn (PBRFT)~\cite{ouyang2022instructgpt}}]
    [{Multi-step / trajectory (Agentic RL)~\cite{zhang2026landscape}}]
  ]
]
\end{forest}
\caption{The proposed three-axis taxonomy for reinforcement learning in
LLM training. The three axes, reward signal (Axis~A), optimization
algorithm (Axis~B), and training scope (Axis~C), are nearly orthogonal,
so any method occupies a single cell of the resulting space.}
\label{fig:taxonomy}
\end{figure*}

\subsection{Axis A: Reward Signal}

The reward is what the policy ultimately optimizes, and its
\emph{source} and \emph{granularity} are the most consequential design
choices. By \emph{source} we distinguish: \emph{human preference}
(RLHF), where annotators rank
responses~\cite{christiano2017deeprl,ouyang2022instructgpt};
\emph{AI preference}, Reinforcement Learning from AI Feedback
(RLAIF) or Constitutional AI, where an LLM judges responses under a
written constitution or rubric~\cite{bai2022constitutional,lee2023rlaif};
\emph{implicit or data preference}, where no reward model is trained and
the preference is encoded directly into the loss, as in
DPO~\cite{rafailov2023dpo}; and \emph{verifiable or rule-based} rewards,
Reinforcement Learning with Verifiable Rewards (RLVR), where correctness
is checked by a deterministic verifier such as a math checker or unit
test~\cite{shao2024deepseekmath,guo2025deepseekr1}.

Orthogonally, by \emph{granularity} we distinguish Outcome Reward
Models (ORMs), which score only the final answer, from Process Reward
Models (PRMs), which score each intermediate
step~\cite{lightman2023verify}. Section~\ref{sec:reward} elaborates on
this axis in detail.

\subsection{Axis B: Optimization Algorithm}

Given a reward signal, the optimization algorithm determines how the
policy is updated. We group methods into three families:
\emph{policy gradient with a critic}
(PPO~\cite{schulman2017ppo}), which estimates advantages with a learned
value function; \emph{critic-free or group-relative} methods
(REINFORCE~\cite{williams1992reinforce}, RLOO~\cite{ahmadian2024rloo},
GRPO~\cite{shao2024deepseekmath}), which avoid a value network by using
baselines or within-group normalization; and \emph{direct or offline}
preference methods (DPO~\cite{rafailov2023dpo}, Identity Preference
Optimization, IPO~\cite{azar2024ipo}, Kahneman--Tversky Optimization,
KTO~\cite{ethayarajh2024kto}), which optimize a closed-form preference
loss without on-policy sampling. The mechanisms and trade-offs of each
family are analyzed in Section~\ref{sec:optimization}.

\subsection{Axis C: Training Scope}

Finally, the training scope fixes the temporal horizon over which the
reward is optimized. \emph{Single-step} methods (PBRFT) treat the task
as one prompt mapped to one response, the degenerate $T{=}1$ MDP
described in Section~\ref{sec:background}. \emph{Multi-step} methods
(Agentic RL) optimize an entire trajectory of interleaved reasoning,
tool calls, and environment actions under a POMDP with $T{>}1$,
introducing planning and long-horizon credit
assignment~\cite{zhang2026landscape}. We dedicate
Section~\ref{sec:agentic} to this emerging paradigm.

\subsection{Placing the Literature}

The power of the taxonomy is that the three axes can be read jointly:
DeepSeek-R1, for instance, is \emph{(A)}~verifiable-reward,
\emph{(B)}~GRPO, \emph{(C)}~single-step (one long reasoning response),
whereas an RL-trained software-engineering agent is
\emph{(A)}~verifiable (unit tests), \emph{(B)}~PPO or GRPO,
\emph{(C)}~multi-step.
Table~\ref{tab:placement} places representative works in this space; the
remaining sections elaborate each row.

\begin{table*}[t]
\centering
\caption{Representative RL-for-LLM methods placed on the three axes.
RLVR = Reinforcement Learning with Verifiable Rewards;
ORM/PRM = Outcome/Process Reward Model;
PBRFT = Preference-Based Reinforcement Fine-Tuning;
CoT = Chain of Thought;
TDPO = Token-level Direct Preference Optimization.}
\label{tab:placement}
\renewcommand{\arraystretch}{1.2}
\footnotesize
\begin{tabular}{@{}llllp{0.30\textwidth}@{}}
\toprule
\textbf{Method} & \textbf{Axis A: Reward} & \textbf{Axis B: Algorithm} & \textbf{Axis C: Scope} & \textbf{Domain / Note} \\
\midrule
InstructGPT~\cite{ouyang2022instructgpt}      & Human pref.\ (ORM)        & PPO            & Single-step & First RLHF assistant at scale \\
Llama~2-Chat~\cite{touvron2023llama2}          & Human pref.\ (ORM)        & PPO + reject.\ sampling & Single-step & Open RLHF at scale \\
Constitutional AI~\cite{bai2022constitutional} & AI pref.\ (ORM)           & PPO            & Single-step & Harmlessness from AI feedback \\
RLAIF~\cite{lee2023rlaif}                       & AI pref.\ (ORM)           & PPO            & Single-step & Matches RLHF on summarization \\
DPO~\cite{rafailov2023dpo}                      & Implicit pref.            & DPO (offline)  & Single-step & No explicit reward model \\
IPO~\cite{azar2024ipo}                          & Implicit pref.            & IPO (offline)  & Single-step & Theoretical fix to DPO overfitting \\
KTO~\cite{ethayarajh2024kto}                    & Implicit (unpaired)       & KTO (offline)  & Single-step & Prospect-theoretic loss \\
TDPO~\cite{zeng2024tdpo}                         & Implicit pref.            & DPO (token)    & Single-step & Token-level control \\
RLOO~\cite{ahmadian2024rloo}                     & Human pref.               & REINFORCE-leave-one-out & Single-step & Critic-free, low variance \\
DeepSeekMath~\cite{shao2024deepseekmath}        & Verifiable (ORM)          & GRPO           & Single-step & Introduces GRPO \\
DeepSeek-R1~\cite{guo2025deepseekr1}            & Verifiable (ORM)          & GRPO           & Single-step$^\dagger$ & Emergent long-CoT reasoning \\
Let's Verify~\cite{lightman2023verify}          & Verifiable (PRM)          & Search / best-of-$n$ & Single-step & Step-level supervision \\
OmegaPRM~\cite{wang2024omegaprm}                 & Verifiable (PRM)          & Auto.\ MC supervision & Single-step & Automated process labels \\
SWE-agent~\cite{yang2024sweagent}               & Verifiable (unit tests)   & PPO / GRPO     & Multi-step  & Agent--computer interface \\
ReAct~\cite{yao2023react}                        & --- (prompting)           & RL-free        & Multi-step  & Reasoning + acting scaffold \\
Toolformer~\cite{schick2023toolformer}          & Self-supervised           & SFT            & Single-step & Learns API calls \\
Voyager~\cite{wang2023voyager}                   & Environment feedback      & RL-free        & Multi-step  & Lifelong embodied agent \\
\bottomrule
\end{tabular}\\[2pt]
{\footnotesize $^\dagger$ A single long reasoning response; the
R1 pipeline also uses multi-stage RL.}
\end{table*}
